% CAL-12: materiał kanoniczny znajduje się w ME-32.
% Wszystkie 23 etykiety oryginalnego rozdziału zachowano w ME-32.
% Do czasu wycofania całej książki pozostawiamy jedynie odsyłacz.
\providecommand{\atlasref}[2]{#2}
\section{Całka Daniella--Stone'a}

Pełną konstrukcję całki Daniella--Stone'a przeniesiono do
\atlasref{ME-32::ch:daniell-stone-32}{Teorii miary, rozdział 32}.
W odróżnieniu od całek uogólniających całkę Riemanna
konstrukcja ta wychodzi od dodatniego funkcjonału liniowego
na kracie funkcji i prowadzi do miary oraz całki Lebesgue'a.
Nie powtarzamy tu definicji, twierdzeń ani dowodów.

W rozdziale docelowym znajdują się:
\begin{itemize}
 \item \atlasref{ME-32::it:def:daniell-functional}{funkcjonał Daniella i ciągłość porządkowa};
 \item \atlasref{ME-32::it:thm:daniell-stone}{konstrukcja i jednoznaczność miary skończonej};
 \item \atlasref{ME-32::it:thm:daniell-unbounded-unit}{wersja dla funkcji nieograniczonych z jedynką};
 \item \atlasref{ME-32::it:thm:daniell-stone-general}{wersja ogólna z warunkiem Stone'a i lokalna jednoznaczność};
 \item \atlasref{ME-32::it:cor:daniell-general-density}{gęstość kraty i uzupełnienie miary};
 \item \atlasref{ME-32::it:thm:daniell-radon-comparison}{porównanie z miarą Radona}.
\end{itemize}
